% TOP_PROBLEM: {"id":30007638,"problem_number":"LOCAL-30007638","title":"Separating discrimination mechanisms","category_id":21,"category":{"id":21,"name":"economics","display_name":"Economics","description":"Open economic theory statements, published research agendas, and proposed scoped specifications; question provenance and historical priority are qualified per record.","slug":"economics","order_index":21,"created_at":"2026-10-08T00:00:00Z"},"status":"research_question","external_url":null,"legacy_ids":[],"published":false,"statement":"In the explicitly specified setting: How can hiring and pay discrimination caused by employer prejudice, customer bias, statistical beliefs, and intersecting identities be separately identified and effectively reduced? The mathematical statement above fixes the target, assumptions and scope of this catalogue formulation.","economics":{"original_id":127,"field":"Labor markets and work","kind":"identification","formalization_basis":"adapted_research_question","source_status":"research_agenda","priority_status":"earliest_found","priority_note":"Earliest directly inspected qualifying passage in this audit; earlier origins were not established. A review or research-agenda statement does not imply original priority.","earliest_verified_reference":{"authors":["Majid Ahmadi","Gwen-Jirō Clochard","Jeff Lachman","John A. List"],"title":"Toward an Understanding of Discrimination When Multiple Channels Exist","year":2025,"url":"https://bfi.uchicago.edu/working-papers/toward-an-understanding-of-discrimination-when-multiple-channels-exist/","doi":"10.3386/w33391","locator":"Author abstract, opening and final paragraphs, dated January 29, 2025","evidence_summary":"Multiple discriminatory channels make source attribution difficult; the paper isolates customer and managerial bias in a particular setting."},"current_reference":{"authors":["Majid Ahmadi","Gwen-Jirō Clochard","Jeff Lachman","John A. List"],"title":"Toward an Understanding of Discrimination When Multiple Channels Exist","year":2025,"url":"https://bfi.uchicago.edu/working-papers/toward-an-understanding-of-discrimination-when-multiple-channels-exist/","doi":"10.3386/w33391","locator":"Author abstract, opening and final paragraphs, dated January 29, 2025","evidence_summary":"Multiple discriminatory channels make source attribution difficult; the paper isolates customer and managerial bias in a particular setting."},"formalization_note":"The cited passage establishes a research direction, not this exact mathematical statement. The notation, population, policy menu, assumptions, and estimands are an original adaptation for this catalogue; no universal unsolved theorem or source priority is claimed."},"statusQualification":"Published question or research agenda verified at the cited locator. The mathematical specification is adapted for this catalogue and includes additional assumptions; source publication does not establish first-ever priority or certify this exact model remains unsolved. The cited passage establishes a research direction, not this exact mathematical statement. The notation, population, policy menu, assumptions, and estimands are an original adaptation for this catalogue; no universal unsolved theorem or source priority is claimed.","statusReviewedAt":"2026-10-08"}
% FIRST_POSTED: 2026-10-08
% ENTRY_KIND: statement_only
% !TeX program = lualatex
\documentclass[11pt]{article}
\usepackage{fontspec}
\usepackage[a4paper,margin=25mm]{geometry}
\usepackage{amsmath,amssymb,mathtools}
\usepackage{xurl}
\usepackage[hidelinks]{hyperref}
\newcommand{\sref}[2]{\hyperlink{source:#1:#2}{[S#2]}}
\setlength{\emergencystretch}{3em}
\begin{document}
% problemId:30007638
\section*{Separating discrimination mechanisms}\label{30007638}
\subsection{Definitions and mathematical statement}
Fix a hiring setting with applicants, employers, and customers. Let \(g_i\) be applicant \(i\)'s recorded identity vector and \(s_i\) independently measured productive characteristics. Let \(H_i(d,c,q)\) indicate hiring under experimentally specified employer information \(d\), customer exposure \(c\), and productivity signal precision \(q\). Employer prejudice concerns identity-dependent preferences holding expected productivity and customer revenue fixed; customer bias concerns the revenue consequences of customer exposure; statistical discrimination concerns decisions based on imperfect productivity beliefs. For identity groups \(g\) and \(g'\), define conditional hiring gaps at common \(s\), and use orthogonal variation in \(d,c,q\) to estimate which channels generate their changes. The identification task is to find assumptions and feasible variations that distinguish these channels, allowing interactions rather than forcing an additive decomposition. A policy arm \(p\) has causal effect \(E[H_i(p)-H_i(0)\mid g_i=g]\), and should also be evaluated for earnings and firm output. Identity labels themselves need not be manipulable. Correspondence tests identify responses to perceived identities under their design; they do not automatically establish motives. Quantify what remains observationally equivalent, especially across intersecting identities and beyond the studied occupation.

\par\textbf{Formulation and scope.} The cited passage establishes a research direction, not this exact mathematical statement. The notation, population, policy menu, assumptions, and estimands are an original adaptation for this catalogue; no universal unsolved theorem or source priority is claimed.

\par\textbf{Open-question provenance.} An explicit question or research agenda was verified at the source locator. This does not by itself certify that every later version remains unresolved \sref{30007638}{1}.

\par\textbf{Historical priority.} Earliest directly inspected qualifying passage in this audit; earlier origins were not established. A review or research-agenda statement does not imply original priority.

\subsection{Short English statement}
In the explicitly specified setting: How can hiring and pay discrimination caused by employer prejudice, customer bias, statistical beliefs, and intersecting identities be separately identified and effectively reduced? The mathematical statement above fixes the target, assumptions and scope of this catalogue formulation.

\subsection{Sources}
\begin{itemize}\raggedright
\item[\textnormal{[S1]}]\hypertarget{source:30007638:1}{} Majid Ahmadi, Gwen-Jirō Clochard, Jeff Lachman, John A. List. 2025. Toward an Understanding of Discrimination When Multiple Channels Exist. Author abstract, opening and final paragraphs, dated January 29, 2025.
\url{https://bfi.uchicago.edu/working-papers/toward-an-understanding-of-discrimination-when-multiple-channels-exist/}.
DOI: 10.3386/w33391.
{\small\textit{Earliest verified question/agenda reference; historical priority qualified below. Multiple discriminatory channels make source attribution difficult; the paper isolates customer and managerial bias in a particular setting.}}
\end{itemize}
\end{document}
